% !PRINTED 68
\weberpage{68}

\begin{center}
{\large \S.\ 20.}

\vspace{0.5em}
{\large\gesperrt{Determinants.}}
\end{center}

We now consider a system of \(n^2\) arbitrary quantities, with which the rational operations of arithmetic can be performed. For a simple notation of these quantities we choose one letter with a double index \(a_i^{(k)}\), where both \(i\) and \(k\) are to run through the series of digits \(1,2,3,\ldots,n\). For better survey we arrange these quantities in a square, so that all \(a\)'s with the same upper index stand in a horizontal row, and all \(a\)'s with the same lower index in a vertical row, and denote this square by \(\Delta\):
\begin{equation}
\Delta=
\begin{vmatrix}
 a_1^{(1)}&a_2^{(1)}&a_3^{(1)}&\cdots&a_n^{(1)}\\
 a_1^{(2)}&a_2^{(2)}&a_3^{(2)}&\cdots&a_n^{(2)}\\
 a_1^{(3)}&a_2^{(3)}&a_3^{(3)}&\cdots&a_n^{(3)}\\
 \vdots&\vdots&\vdots&&\vdots\\
 a_1^{(n)}&a_2^{(n)}&a_3^{(n)}&\cdots&a_n^{(n)}
\end{vmatrix}.
\tag{1}
\end{equation}
For brevity one calls the horizontal rows \gesperrt{rows}, and the vertical rows \gesperrt{columns}.

But by the square enclosed between vertical lines we shall understand not only the complex of the quantities \(a\), but a definite arithmetical combination of these quantities, which can be calculated as soon as the \(a\)'s are given numerically, and which we shall now describe.

Form the product of the terms standing in the diagonal running from upper left to lower right:
\begin{equation}
 M=a_1^{(1)}a_2^{(2)}a_3^{(3)}\cdots a_n^{(n)}.
\tag{2}
\end{equation}
From this derive \(\Pi(n)\) products \(M,M',M'',\ldots\) by permuting the lower indices, and give to each product so obtained the \gesperrt{positive or negative sign} according as the applied permutation belongs to the first or to the second kind; thus, in the notation of the preceding paragraph,
\begin{equation}
 M'=\pm a_{\alpha_1}^{(1)}a_{\alpha_2}^{(2)}\cdots a_{\alpha_n}^{(n)}.
\tag{3}
\end{equation}
The sum of these products
\[
M+M'+M''+\cdots=\Sigma M
\]
%CONT>>
